% !TEX program = xelatex
% Single-side, single-page A4 summary for 良文杯 2026 finals
% Team 792 "超级回天什么时候上映" — 史景喆，清华大学
% Content distilled from report_full.tex (Features / Model / Results only)
\documentclass[10pt,twocolumn]{article}

\usepackage[a4paper,left=0.40in,right=0.40in,top=0.40in,bottom=0.40in,columnsep=0.18in]{geometry}

\usepackage{fontspec}
\usepackage{xeCJK}
\setCJKmainfont[BoldFont={Noto Serif CJK SC Bold},ItalicFont={Noto Serif CJK SC}]{Noto Serif CJK SC}
\setCJKsansfont{Noto Sans CJK SC}
\setCJKmonofont{Noto Sans Mono CJK SC}
\setmainfont{Nimbus Roman}[
  Extension={.otf},
  Path=/usr/share/fonts/opentype/urw-base35/,
  UprightFont=NimbusRoman-Regular,
  BoldFont=NimbusRoman-Bold,
  ItalicFont=NimbusRoman-Italic,
  BoldItalicFont=NimbusRoman-BoldItalic,
]

\usepackage[T1]{fontenc}
%   
\usepackage{xcolor}
\usepackage{titlesec}
\usepackage{booktabs}
\usepackage{array}
\usepackage{enumitem}
\usepackage{caption}
\usepackage{tabularx}
\usepackage{pifont}
\usepackage{amsmath}
\usepackage{amssymb}
\usepackage{url}

% Tighten everything
\setlength{\parindent}{0pt}
\setlength{\parskip}{1pt}
\setlength{\abovedisplayskip}{1pt}
\setlength{\belowdisplayskip}{1pt}
\setlength{\abovedisplayshortskip}{0pt}
\setlength{\belowdisplayshortskip}{0pt}
\setlength{\textfloatsep}{2pt}
\setlength{\intextsep}{2pt}
\setlength{\floatsep}{2pt}
\captionsetup{font=scriptsize,skip=1pt,labelfont=bf}
\setlist[itemize]{leftmargin=*,topsep=0pt,itemsep=0pt,parsep=0pt,labelsep=4pt}

% ICML-like section headings
\definecolor{icmlblue}{RGB}{0,30,90}
\definecolor{accred}{RGB}{170,0,0}
\titleformat{\section}{\normalfont\bfseries\color{icmlblue}\small}{\thesection.}{0.3em}{}
\titlespacing*{\section}{0pt}{3pt}{1pt}

\pagestyle{empty}

\begin{document}

\twocolumn[{%
\centering
{\large\bfseries\color{icmlblue}面向 PnL 目标的 LOB 中间价方向预测：$\Delta$mid 回归 $+$ SPO+ 决策焦点学习 $+$ $50\!\times\!50$ 异质集成}\par
\vspace{1pt}
{\footnotesize\itshape 第二届"良文杯"统计建模与 AI 预测挑战赛\,\textemdash\,汇报摘要}\par
\vspace{2pt}
{\scriptsize 随机码 \textbf{792} \,\textbar\, 队伍：\textbf{超级回天什么时候上映} \,\textbar\, 作者：\textbf{史景喆}（清华大学）\,\textbar\, 公榜 $-6.65 \to$ \textcolor{accred}{\textbf{$+35.64$}}（增益 \textcolor{accred}{\textbf{$+42.29$}}），\textbf{公榜 \#1}，私榜 \textbf{\#1}（待组委会公布分数）}\par
\vspace{3pt}
}]

\scriptsize

\section{摘要}

我们将赛题理解为：在少量股票样本上训练、在其他股票上预测的、带有交易成本的期望收益预测问题；研究过程紧盯\textbf{正确性、有效性、范化性}。最终方案以因子工程 + L2 $\Delta$mid 回归 + SPO+ 决策焦点学习 (DFL) + 50 NN $\times$ 50 LightGBM 异质集成构成，\textbf{公榜累计收益率 $+35.6$（baseline $-6.7$），公榜与私榜均为决赛参赛队伍第一}。

我们同时基于 Claude Code 的 Agent Harness 构建了 Manager-Worker 多 Agent 架构、并集成至飞书 IM，实现\textbf{正确性、高效性、可追溯性}的 Coding Agent 研究流程：一周内完成约 200 组本地实验、约 200 次 commit。代码开源于 \url{https://github.com/JingzheShi/liangwenbei}。




% =========================================================
\section{Problem}

数据：5 只匿名 A 股 $\times$ 120 交易日 $\times$ 上午/下午各 2001 ticks（3s/tick），共 1200 parquet（约 474 MB），154 维原始字段（OHLC / 10 档 LOB 报价与挂单量 / 6 类订单流强度・指示・加速度 / 衍生 spread・imbalance 等）；1192 文件无 NaN、8 文件涨停日深档 ask 含少量缺失（pipeline 用 \texttt{isfinite} 兜底）。
评分：单笔 $\mathrm{pnl}_i^{(h)}=\frac{(\hat a_i-1)\Delta\mathrm{mp}_h^{(i)} - 10^{-4}|\hat a_i-1|((\mathrm{mp}_{t+h}{+}1)+(\mathrm{mp}_t{+}1))}{\mathrm{mp}_t+1}$，$\hat a\in\{0,1,2\}$；买/卖才扣 $\sim 0.02\%$ 双边手续费。总分 $\mathrm{Score}=\max_{h\in\{5,10,20,40,60\}}\sum_i \mathrm{pnl}_i^{(h)}$，5 horizon 取最大——本文以 $h=60$ 为主线。


% =========================================================
\section{因子工程}
\paragraph{六大家族，359 维。}以 154 维原始字段为底，叠加 216 维派生量、按 KS 检验丢弃 11 个分布偏移敏感列，得到 SchemeP 359 维：\textbf{LOB 派生}（$\sim$30 维：midprice / spread / imbalance / 微价格）、\textbf{多尺度 OFI}（$\sim$50 维：MLOFI / EWMA-OFI 多 $\alpha$ / Kyle $\lambda$ rolling，$W\!\in\!\{5,20,50,100\}$）、\textbf{订单流强度}（$\sim$30 维：6 类订单 \texttt{intst/ind/acc} + cancel pressure）、\textbf{微结构波动}（$\sim$25 维：signed RV / bipower / Roll / vol-burst）、\textbf{窗口统计}（$\sim$60 维：100-tick z-score / dual-z / quantile rank）、\textbf{不对称性}（$\sim$30 维：bid/ask asymmetry / GOFI / EWMA residual）。

\paragraph{窗口内 z-score (stock-agnostic)。}所有统计量在 100-tick 切片内做自适应 z-score——跨 stock 分布差异被切片内归一化吸收，避免任何 per-stock 统计量。我们尝试过全局 z-score（全训练集 mean/std），在跨 stock 留一验证上明显更差。

\paragraph{方向对称增强 (aug\_a)。}训练时把每个样本与其"bid/ask 互换且 label 反转"的镜像样本拼接，训练集扩大 $2\times$。物理上 LOB 在 bid/ask 互换且标签反向后对应同一市场状态——强制学到方向对称表示，显著降低单边偏差。

\paragraph{重尾压制。}\texttt{amount\_delta}（唯一未官方归一化字段，量级 $10^3$--$10^6$）在 \texttt{raw\_last} 处做保号 log：$v\mapsto\operatorname{sign}(v)\log(1+|v|)$。

% =========================================================
\section{模型工程}

最终模型由 LightGBM 与 MLP 两类基础模型构成，族内 simple mean，跨族固定权重。

\paragraph{NN：4-layer MLP。}$359 \to 256 \to 128 \to 64 \to 1$，LayerNorm + GELU + Dropout（$\sim$120k 参数）。LayerNorm 在 sample 级归一化，杜绝 batch 间数据污染。我们刻意不使用 RNN/Transformer/CNN——窗口内因子已蒸馏 100-tick 时序信息，更复杂架构反带来 in-sample 拟合膨胀。这与近期 tabular DL 文献一致：适当训练下纯 MLP 在 LOB / 表格上即可达到 SOTA。

\paragraph{LGB：L2 $\Delta$mid 回归。}\texttt{objective=regression\_l2}，目标
\begin{equation}\label{eq:label}
y = \frac{\mathrm{mp}_{t+60} - \mathrm{mp}_t}{\mathrm{mp}_t + 1}.
\end{equation}
以 h=60 三分类标签频率倒数为样本权重（class-balanced），使稀疏的涨/跌样本获更高梯度。


\section{训练和模型集成}

\paragraph{树模型：$\Delta$mid L2 回归。}相比 3-class CE，L2 回归保留 $\Delta$mid 的幅度信息，决策层能权衡交易成本与潜在收益；单 LGB 单 seed 下 CE 反而略胜，但在 50$\times$50 集成内 L2 提供与 NN+SPO+ 兼容的连续 $\Delta$mid 信号，决策层能权衡幅度，最终公榜 $+35.64$。

\paragraph{神经网络：SPO+ 决策焦点学习 (DFL)。}L2 回归梯度未直接对齐离散决策。两阶段微调：Phase 1 以 L2 预训练 MLP；Phase 2 warm-start 联合 L2 + SPO+ surrogate。给定单边手续费率 $f$，oracle $z^{\star}(y)=\operatorname{sign}(y)\,\mathbf{1}[|y|>f]$；SPO+ surrogate（1D 三动作改写）与联合损失：
\[
\ell_{\mathrm{SPO+}}(\hat y,y;f) = \mathrm{ReLU}(|2\hat y{-}y|-f) - z^{\star}(y)(2\hat y{-}y) + f|z^{\star}(y)|,
\]
\[
L = (\hat y-y)^2 + \lambda\,\ell_{\mathrm{SPO+}},\quad \lambda=30,\ \mathrm{lr}=3{\times}10^{-5},\ 11\,\mathrm{ep}.
\]

\paragraph{全量数据 retrain。}调参/早停阶段使用 \texttt{train 0--79 / val 80--95}；最终提交将训练区推到 \texttt{date 0--119} 全量，iteration 数取部分训练的 $1.1\times$。

\paragraph{集成。}5 组超参 $\times$ 10 seed 训练 50 LGB 与 50 MLP，跨族固定 $\mathrm{LGB:NN}=3:2$ 的 simple mean。\emph{不在 OOF 上学决策参数与权重}：NN-NN avg corr $0.93$、LGB-LGB $0.84$、设计矩阵条件数 $5.4{\times}10^{4}$，bounded NNLS 学权重 in-sample $+180/+194$ 但公榜 $-4.64$——病态共线性下 simple mean 是数学上的 OOD 稳健最优。




% =========================================================
\section{执行侧}

集成预测 $\hat y$ 通过非对称 EV gate 解码为离散动作：
\[
\hat a =
\begin{cases}
+1\,(\text{买}), & \hat y > +\theta_{\mathrm{up}},\\[-1pt]
-1\,(\text{卖}), & \hat y < -\theta_{\mathrm{dn}},\\[-1pt]
\phantom{+}0\,(\text{不出手}), & \text{otherwise.}
\end{cases}
\]
$\theta_{\mathrm{up}}=3.00{\times}10^{-4}$、$\theta_{\mathrm{dn}}=2.16{\times}10^{-4}$ 由验证集 2D DE 搜出后冻结。两者不对称（涨方阈值更高）以匹配训练样本的正向漂移偏差。详见 \S 实验结果 大 ablation 表 ensemble 块内 sym vs asym EV 的本地 V4 对照。

% =========================================================
\section{实验结果}

\paragraph{大 ablation 表（V4 walk-forward，LGB vs NN 渐进特征 $+$ Ensemble）。}train 0--79 / val 80--95（ES、阈值搜索）/ test 96--119（test 零调参）；指标 $\mathrm{Score}^{(h=60)}_{\text{test}}=\sum_i\mathrm{pnl}_i^{(60)}$，0.02\% 双边手续费。\textbf{L1--L8 报 seed 1--5 mean $\pm$ std}（其余 HP 一致；LGB 用 \texttt{HP\_CONFIGS[0]}, L2 obj, sym EV; NN 4-layer MLP, L2 obj + window-$z$, sym EV）；ensemble 行本身为 multi-seed 聚合。

\begin{center}
\setlength{\tabcolsep}{3pt}\renewcommand{\arraystretch}{0.95}
\scriptsize
\begin{tabularx}{\linewidth}{@{}c >{\raggedright\arraybackslash}X c r r@{}}
\toprule
Lvl & Feature set & dim & LGB Test PnL & NN Test PnL \\
\midrule
\multicolumn{5}{@{}l}{\textit{\color{icmlblue}Progressive features — 154 raw $\to$ 6 family 叠加 $\to$ 370 full $\to$ 359 final（同 L2 head, sym EV）}} \\
L1 & 154 raw                                       & 154 & $+17.70{\scriptscriptstyle\,\pm 1.33}$  & $+26.05{\scriptscriptstyle\,\pm 0.57}$ \\
L2 & $+$ LOB 派生 (midprice / spread / imbalance)   & 165 & $+17.53{\scriptscriptstyle\,\pm 1.40}$  & $+25.90{\scriptscriptstyle\,\pm 0.37}$ \\
L3 & $+$ 多尺度 OFI (MLOFI / EWMA-OFI / Kyle $\lambda$) & 213 & $+22.61{\scriptscriptstyle\,\pm 0.61}$  & $+27.53{\scriptscriptstyle\,\pm 1.22}$ \\
L4 & $+$ 订单流强度 (intst / cancel pressure)        & 240 & $+25.72{\scriptscriptstyle\,\pm 1.89}$  & $+29.86{\scriptscriptstyle\,\pm 0.88}$ \\
L5 & $+$ 微结构波动 (signed RV / vol-burst / Roll)   & 272 & $+24.93{\scriptscriptstyle\,\pm 1.42}$  & $+28.73{\scriptscriptstyle\,\pm 1.63}$ \\
L6 & $+$ 窗口统计 (100-tick z / dual-z / qrank)      & 335 & $+26.31{\scriptscriptstyle\,\pm 0.38}$  & $+31.38{\scriptscriptstyle\,\pm 1.31}$ \\
L7 & $+$ 不对称性 (GOFI / liq-asym) $=$ 370 full     & 370 & $+27.28{\scriptscriptstyle\,\pm 2.34}$  & $+34.77{\scriptscriptstyle\,\pm 1.39}$ \\
L8 & $-11$ KS-fail $=$ 359 final SchemeP            & 359 & $+26.58{\scriptscriptstyle\,\pm 0.96}$  & $+34.34{\scriptscriptstyle\,\pm 0.50}$ \\
\midrule
\multicolumn{5}{@{}l}{\textit{\color{icmlblue}Ensemble — NN $\times$ LGB 跨族（359-d SchemeP, V4 test）}} \\
\multicolumn{3}{@{}l}{(1$+$1) NN $+$ LGB, sym EV}                              & \multicolumn{2}{r}{$+35.02$} \\
\multicolumn{3}{@{}l}{(5$+$5) NN $+$ LGB, sym EV}                              & \multicolumn{2}{r}{$+36.07$} \\
\multicolumn{3}{@{}l}{(5$+$5) NN $+$ LGB, asym EV}                             & \multicolumn{2}{r}{$+36.19$} \\
\multicolumn{3}{@{}l}{(5$+$5) NN$+$SPO $+$ LGB, asym EV}                       & \multicolumn{2}{r}{$\mathbf{+37.73}$ (SOTA)} \\
\bottomrule
\end{tabularx}
\captionof{table}{\textbf{大 ablation：L1--L8 渐进特征 LGB vs NN 双柱对照（5-seed mean $\pm$ std）+ ensemble.} 同一 feature partition、同一 V4 split、同 L2 obj + sym EV 阈值搜，仅 backbone 不同。\textbf{LGB 与 NN 在每档都同序增长}（154$\to$370），但 NN $+8.3$ 高出 LGB——树模型对线性可分的 micro-structure 信号不如 MLP。\textbf{NN window-$z$ critical}：去掉 window-$z$，NN L1 (154 raw + L2 sym EV) 从 $+26.05$ 崩到 $+0.98{\scriptscriptstyle\,\pm 0.20}$（跃升 $+25.07$）；首层 Linear 在 scale 失配输入上无法恢复，LayerNorm 在其后；树则 axis-aligned scale-invariant、无此现象（同步骤 $\pm 0.4$ within noise）。L8 = L7 $-$ 11 KS-fail，NN/LGB 均 $\pm 0.5$ within noise（KS drop 旨在缓解 OOD，本地无明显损失）。Ensemble：(1$+$1) 跨族即达 $+35.02$；扩到 $5{+}5$ multi-seed $+1.05$；asym EV $+0.12$；SPO 决策焦点 $+1.54$，累计 $\mathbf{+37.73}$（SOTA）。LGB 训练用项目 \texttt{HP\_CONFIGS[0]}（GPU）。}
\end{center}

\begin{center}
\setlength{\tabcolsep}{3pt}\renewcommand{\arraystretch}{1.05}
\begin{tabularx}{\linewidth}{@{}l c r r >{\raggedright\arraybackslash}X@{}}
\toprule
Board & $h$ & Cum.\ PnL & $\alpha$/trade & Notes \\
\midrule
Public  & $60$ & $35.64$ & $1.91{\times}10^{-4}$ & \#1（公榜最高）\\
Private & $60$ & $41.61$ & $2.39{\times}10^{-4}$ & \#1（私榜整体最高）\\
Private & $40$ & $38.65$ & $1.94{\times}10^{-4}$ & \#1（私榜 $h{=}40$ 最高）\\
\bottomrule
\end{tabularx}
\captionof{table}{\textbf{Leaderboard.} 本队（792 - 超级回天什么时候上映）在公榜（$h{=}60$）与私榜（$h{=}60$、$h{=}40$）均位列第一。}
\end{center}

\paragraph{多 horizon 阈值缩放。}5 个 horizon 共享同一套 $h{=}60$ 回归预测；阈值 $(\theta_{\mathrm{up}}, \theta_{\mathrm{dn}})$ 按 $\sqrt{h/60}$ 缩放到 $h\!\in\!\{5,10,20,40\}$。这使得 $h{=}40$ 也在私榜上拿到第一（$+38.65$），而无需为短 horizon 单独训练模型。

\section{Coding Agent 在量化中的使用}

\paragraph{系统设计。}基于 Claude Code Agent Harness 构建 Manager-Worker 双层架构并集成至飞书 IM：(i) IM 直连，手机/电脑通过飞书与 Agent 交互；(ii) Manager-Worker 分层，Worker 细节不污染 Manager 上下文；(iii) GitHub + WandB 版本/可视化管理。

\paragraph{经验。}(1) Agent 对\emph{宽且浅}的工作高效（批量 ablation / 归并）；(2) \emph{窄且深}的研究决策仍依赖人类先验（"决策参数不在 OOF 上学"、"公榜是唯一 OOD 裁判"）；(3) \textbf{硬约束须以"协议"形式注入每个 worker prompt}——\texttt{CRITICAL\_CONSTRAINTS.md} 三条若未注入，worker 可能在不知情时违反平台协议。一周内完成约 200 组本地实验、约 200 次 commit，公榜从 $-6.7$ 提升到 $+35.6$。



\end{document}

